% Addendum to FF.45; append after its proof and before the source-scope paragraph.
\paragraph{The remaining action is trivial on all constant-base \(p\)-Kummer groups.}
Let \(E/B_p\) be finite and let \(\varphi\) be in the
pointwise kernel of Claim~\ref{clm:FFgenerickernel} for
\(\mathbb A^1_E\). For every finite \(L/E\), every
positive integer \(n\), and every \(r\in L(T)^\times\),
its generic action satisfies
\[
       \alpha_\varphi^*\kappa_{p^n}(r)
                         =\kappa_{p^n}(r)
              \quad\hbox{in }H^1(G_{L(T)},\mu_{p^n}).
                                                        \tag{FF.45.8}
\]
In particular, the action is the identity on the maximal
abelian pro-\(p\) quotient after every finite constant
base change. This conclusion concerns abelian quotients,
not the entire geometrically pro-\(p\) group.

Here is the proof, including descent from split divisors.
For \(a\in L\), pointwise identity makes the inverse
image of \(\mathbb A^1_L\setminus\{a\}\) itself.
Its geometric integral \(p\)-adic first cohomology has
rank one. The argument giving (FF.44.5) therefore yields
\[
 \varphi^*\kappa_p(T-a)
       =u_a\kappa_p(T-a)+\eta_a,
 \qquad u_a\in\mathbb Z_p^\times,
 \quad\eta_a\in H^1(G_L,\mathbb Z_p(1)).
                                                        \tag{FF.45.9}
\]
Evaluation at \(T=a+1\), whose actual value is unchanged,
gives \(\eta_a=0\). Choose \(z\in L^\times\setminus
W_p\). Such an element exists because \(G_L\) is a
nontrivial free pro-\(p\) group, hence has a nontrivial
\(\mathbb F_p\)-quotient, and Kummer theory applies.
By Claim~\ref{clm:FFmetabeliankernel},
\(\kappa_p(z)\ne0\). Evaluation at \(a+z\) gives
\((u_a-1)\kappa_p(z)=0\). The group
\(H^1(G_L,\mathbb Z_p(1))\) is
\(\operatorname{Hom}_{\rm cont}(G_L,\mathbb Z_p(1))\),
because \(\mu_{p^\infty}\subset L\), and is
torsion-free over \(\mathbb Z_p\). Thus \(u_a=1\).
This proves that every linear-unit integral Kummer class
is fixed, and hence that all of its finite-level classes
are fixed.

For general \(r\), let \(V\subset\mathbb A^1_L\)
be its nonvanishing locus and choose a finite Galois
extension \(L'/L\) splitting its zero and pole divisor.
On \(V_{L'}\) it is a constant times a product of
integer powers of the linear units already considered.
The relative structural isomorphism fixes every constant
Kummer class. Therefore
\[
 \delta=\varphi^*\kappa_{p^n}(r)-\kappa_{p^n}(r)
                          \in H^1(V,\mu_{p^n})
\]
restricts to zero on \(V_{L'}\). The
inflation--restriction sequence for this purely constant
Galois cover puts \(\delta\) in the image of
\(H^1(\operatorname{Gal}(L'/L),\mu_{p^n})\), hence
in the constant cohomology subgroup from \(G_L\).
Choose \(t_0\in V(L)\), possible since \(L\) is
infinite. Its actual point is fixed, so
\(t_0^*\delta=0\). Evaluation is a left inverse to
inflation of constant classes; consequently
\(\delta=0\). Restriction to the generic point proves
(FF.45.8). Finally, every class in
\(H^1(G_{L(T)},\mu_{p^n})\) is a Kummer class, and
\(\mu_{p^n}\subset L\), so (FF.45.8) is exactly the
asserted identity on finite abelian \(p\)-quotients.

One can also identify the geometric inertia exponents
before specializing to \(p\). On a projective line with
finitely many geometric punctures, the inertia classes
generate geometric abelianization with the single
relation that their sum is zero. An automorphism fixing
the punctures and multiplying their generators by units
must therefore use the same unit at every puncture.
Compatibility when punctures are added makes these
units one element of \(\widehat{\mathbb Z}^{\times}\).
The preceding calculation proves that its \(p\)-component
is one. It does not determine its other components.

The distinction from inner triviality is concrete even
for an abstract four-punctured line. Let \(P\) be the
free pro-\(p\) group on \(x,y,z\), with fourth
peripheral generator \((xyz)^{-1}\). The automorphism
\[
 x\longmapsto(xy)x(xy)^{-1},\qquad
 y\longmapsto(xy)y(xy)^{-1},\qquad
 z\longmapsto z
                                                        \tag{FF.45.10}
\]
has inverse given by conjugating the first two generators
by \((xy)^{-1}\). It fixes \(xy\), \(xyz\), every
peripheral conjugacy class, and abelianization. It is not
inner. In the torsion-free class-two pro-\(p\) quotient,
the independent central commutators are
\([x,y],[x,z],[y,z]\). Conjugation by an element with
abelian coordinates \((a,b,c)\) could agree with
(FF.45.10) on \(x,y\) only if \(a=b=1,c=0\).
Its commutator with \(z\) would then be
\([x,z][y,z]\ne1\), a contradiction. This quotient
exists also for \(p=2\): use coordinates in
\(\mathbb Z_p^3\times\mathbb Z_p^3\), with the
central product correction
\((a_i a'_j)_{i<j}\), so no division by two is involved.
This is an abstract obstruction to an abelianization-only
argument. No compatibility of (FF.45.10) with the actual
\(G_E\)-action, with every finite cover, or with the
full affine-line etale topos is asserted.

The class-two calculation was executed modulo 2, 3, 4, 5 and 8. The order-64 model was checked on all 4096 products. The program is

\path{agent_F/F_metabelian_checks.py}.
